\begin{abstract}

    %This Systematic Mapping Study (SMS) explores the importance of serious and educational digital games in raising awareness, preventing, and combating sexual violence. The main objective of the research is to identify and characterize existing games in this context, evaluating their effectiveness, impact, and development approaches. Through rigorous inclusion and exclusion criteria, the study seeks to answer key questions about narrative styles, formats, and platforms of the games, aiming to understand how these elements affect user experience and the conveyance of educational messages. With an in-depth analysis of the selected articles, this work intends to provide valuable insights for researchers, professionals, educators, and game developers, contributing to the advancement of knowledge in the area and guiding future research and practices.
    
    Generative Artificial Intelligence (GenAI) has established itself as a major transformative vector in non-interactive audiovisual production, driven by the technical maturation of architectures such as Diffusion Models and Transformers. Although the adoption of these tools in real production workflows has grown significantly, the literature remains fragmented, and existing taxonomic proposals address technical capabilities, production \textit{pipeline} stages, and ethical risks in isolation, which prevents any joint analysis of the \textit{trade-offs} entailed by each technological choice. In light of this scenario, this research develops TIGA (Taxonomy for Integrating Generative AI in Non-Interactive Audiovisual Production), a multidimensional taxonomy built through the iterative taxonomy development method upon the empirical basis of a previous Systematic Mapping Study (SMS) that consolidated sixty primary studies retrieved within a 2022 to April 2025 search window and actually published between 2023 and April 2025. Given the growing volume of publications and the rapid obsolescence of commercial tools, methodological rigor was concentrated on the qualitative synthesis and the analytical reinterpretation of the previously extracted data, subordinated to the meta-characteristic \textit{integration of GenAI into the non-interactive audiovisual production pipeline}, rather than on new quantitative search rounds. TIGA is faceted: each instance simultaneously receives an operational profile (degree of autonomy, \textit{pipeline} stage, and interaction mode), a technical profile (architectural family, base model, output type, infrastructure, and dominant bottleneck), and a risk and ethics profile (authorship, bias, impact on labor, transparency, reproducibility, and perceptibility), complemented by a decision rubric, an application protocol, and a \textit{trade-off} map. Applying the taxonomy to three real studies from the corpus confirmed the correlation between greater autonomy delegated to the machine and higher sociotechnical cost. This document delivers the preliminary version of the artifact and the instruments for its evaluation; full validation, through an expert panel and an inter-rater agreement measure, constitutes the next stage of the research.
    
    
    
    \end{abstract}
    